\documentclass{umk-matematika}
\begin{document}

\begin{center}
{\Large\bfseries Практическая работа 14}
\end{center}
\vspace{0.5cm}

\textbf{Задача 1.} Что такое двугранный угол?\par\vspace{10pt}
\textbf{Задача 2.} Что такое линейный угол двугранного угла?\par\vspace{10pt}
\textbf{Задача 3.} Чему равен двугранный угол между стеной и полом в комнате?\par\vspace{10pt}
\textbf{Задача 4.} Какие плоскости называются перпендикулярными?\par\vspace{10pt}
\textbf{Задача 5.} В кубе $ABCDA\_1B\_1C\_1D\_1$ найдите линейный угол двугранного угла между плоскостями ABC и AB$_1$C$_1$.\par\vspace{10pt}
\textbf{Задача 6.} Двугранный угол равен 60$^\circ$. На одной грани взята точка A, удалённая от ребра на 10 см. Найдите расстояние от точки A до другой грани.\par\vspace{10pt}
\textbf{Задача 7.} Признак перпендикулярности двух плоскостей.\par\vspace{10pt}
\textbf{Задача 8.} В правильной треугольной призме сторона основания равна 4 см, боковое ребро — 6 см. Найдите двугранный угол между плоскостью основания и плоскостью, проходящей через сторону основания и середину противолежащего бокового ребра.\par\vspace{10pt}
\textbf{Задача 9.} Двугранный угол равен 120$^\circ$. Точка A лежит внутри угла и удалена от обеих граней на расстояние 4 см. Найдите расстояние от точки A до ребра двугранного угла.\par\vspace{10pt}
\textbf{Задача 10.} Крыша имеет форму двугранного угла 120$^\circ$. На скате крыши на расстоянии 3 м от конька (ребра) установлена опора, перпендикулярная скату. На какой высоте от основания крыши находится точка крепления опоры, если ширина дома 8 м?\par\vspace{10pt}
\newpage
\section*{Ответы}

\begin{enumerate}
  \item Ответ: две полуплоскости с общей границей.
  \item Ответ: угол между перпендикулярами к ребру.
  \item Ответ: 90$^\circ$.
  \item Ответ: двугранный угол = 90$^\circ$.
  \item Ответ: 90$^\circ$.
  \item Ответ: $5\sqrt{3}$ см.
  \item Ответ: признак перпендикулярности плоскостей.
  \item Ответ: $60^\circ$.
  \item Ответ: 8 см.
  \item Ответ: $\approx$ 4,9 м.
\end{enumerate}
\end{document}
